\documentclass[11pt]{report}
% ============================================================
%  PREAMBOLO VERIFICA — font normale (TeX Gyre Pagella)
% ============================================================
\usepackage{fontspec}
\setmainfont{TeX Gyre Pagella}

\input{../../preambolo-verifica-core.tex}

% unicode-math va caricato dopo amsmath/amssymb/amsfonts (già caricati
% dal core), altrimenti conflitti tipo "\eth already defined".
\usepackage{unicode-math}
\setmathfont{TeX Gyre Pagella Math}


\begin{document}
\chapter{Valutazioni Interrogazioni di Recupero}

\begin{description}
    \item[CLAMAR] Buono. Sorridente e positiva. \textbf{Proposta: 8-}
    
    \begin{enumerate}
    \item $-\frac{1234}{5678}x^2-2=0$
    Ha fatto scomparire il coeff di x² e ha inizialmente messo la radice facendo scomparire il meno.

    Quando le ho suggerito di portare di là il 2 ci è arrivata (le ho detto di fare attenzione ai segni)
    \item $-2(x+2)^2=-8+x$ Errorino di segno ma solo copiando. Ottimo.
    \item $4+(4x-2)(4x+2)=0$. Si ricordava il prodotto notevole. Perfetto senza aiuti.
    \item $2x^2-\pi=0$. Perfetto senza aiuti.
    \end{enumerate}
    \item[GUGGOL] \textbf{Proposta: 8-}
    
    \begin{enumerate}
        \item $(2x-1)^2=(\frac 12x+1)^2$. Ha sbagliato tutti i quadrati dei binomi. Ha fatto le somme e sottrazioni tra frazioni con la virgola. Poi però gli ho chiesto di farlo in frazione e ci è riuscito. In realtà tutto ciò che riguarda un'eq spuria lo ha fatto giusto.
        \item $\sqrt 2 x^2 - 2 =0$. Ha ragionato bene, abbiamo cosmetizzato il $2/{\sqrt 2}.$ Lui ha voluto prendere il quadrato da entrambe le parti per far scomparire la radice così ha ottenuto un'equazione quarta. gli ho chiesto se dovesse mettere il $\pm$ e mi ha detto sì perché l'indice è pari.
        \item $(2x+1)(\frac 14+x)=x(\frac 32-3x).$ Ha sbagliato a copiare una cosuccia ma solo distrazione. è arrivato da solo a \dots
        \item[CARANT] \textbf{Proposta: }
        \begin{enumerate}
            \item $x^2+x=2x$ Ha portato $2x$ dall'altra parte e ha fatto tutto da sola perfettamente.
            \item $\frac{x^2}{2}+2=1$ Non sapeva come occuparsi del 2 a denominatoere. Non le piacciono le frazioni, non si ricordava come sommare le frazioni. Alla fine abbiamo deciso di moltiplicare tutto per $2$. Si è ricordata che x²=-2 era impossibile.
        \end{enumerate}
        \item[GUICER] \textbf{Proposta: 8}
        \item[RICRAO] \textbf{Proposta: 6}
        \item[VALCAS] \textbf{Proposta: 6}  
    \end{enumerate}
\end{description}

\end{document}